\subsection{Explicit inverse moduli for the two charts}
The same rectangular parameter $(s,\theta)$ represents the spatial and
planar quarters, with $v=\sin\theta$, $w=\cos\theta$, $d=1+3v^2$.
Suppose $\theta\le\eta$ and put $\delta=\eta-\theta$. For arbitrary
$r,s\in[0,1]$ we claim
\begin{equation}
 |X(r,\eta)-X(s,\theta)|\ge \delta^4/2048,\qquad
 |F(r,\eta)-F(s,\theta)|\ge \delta^5/4096 .
 \label{eq:invmod}
\end{equation}
These deliberately coarse constants include the two singular corners.

Take the spatial supporting normal $n(\theta)=(-c,\sin t,\sqrt{1+2c})$.
Its projected lengths on the two circle planes are $1,1+c$, and
$|n|=\sqrt{2(1+c)}\le2$. The support deficiency of the point
$X(r,\eta)$, measured from the plane containing $X(s,\theta)$, is
\[
 \frac{1-r}{2}|A(\eta)-A(\theta)|^2
     +\frac{r(1+c)}{2}|B(\eta)-B(\theta)|^2 .
\]
Indeed a unit circle point at chord distance $b$ from its maximizer
loses $b^2/2$ in unit support height. Here $A(\xi), B(\xi)$ denote the
two spatial endpoints at parameter $\xi$. The absolute derivatives of
the two angular arclengths are respectively $2\sqrt3 w/d$ and
$2\sqrt3 v/\sqrt d$; the second follows from
$\theta^*=\pi/2-\beta$, $d^*=4/d$. They are bounded below by
$(\sqrt3/2)\cos\theta$ and $\sqrt3\sin\theta$. The integrals of either
sine or cosine on any subinterval of length $\delta$ of $[0,\pi/2]$
are at least $\delta^2/4$, since $\sin x\ge x/2$ there. Both circular
angle increments exceed or equal $\delta^2/8$, and both chord lengths
are at least $\delta^2/16$. The inequalities $1+c\ge1/2$ and $|n|\le2$
now yield the first bound.

For the plane write $e=(e_U,e_V)$ for the unit ruling direction in
$(U,V)$ coordinates. The perpendicular $\nu=(e_V,-e_U)$ has
nonnegative entries and $\nu_U\ge1/4$ from
$e_V=(1+v^2+2v^4)/d$. Both coordinates of both endpoints $a,b$
increase with $\theta$. Their $U$ derivatives are
\[
    a_U'=2\sqrt3 w^4/d^2\ge(\sqrt3/8)w^4,\qquad
    b_U'=4\sqrt3 v^4/d\ge\sqrt3 v^4 .
\]
Now take the scalar product of
$F(r,\eta)-F(s,\theta)$ with $\nu(\theta)$; the contribution of
$F(r,\theta)-F(s,\theta)$ vanishes. The remaining projection is at least
a quarter of the $U$ increment, a convex combination of the two endpoint
increments. On any interval of length $\delta$, the integrals of
$\sin^4\theta$ and $\cos^4\theta$ are each at least $\delta^5/80$,
using $\sin x\ge x/2$ again. This proves the planar inequality.
For either $Y=X,F$ we also have
\[
 \sqrt3|s-r|\le |Y(s,\theta)-Y(r,\eta)|+
                       4\delta .
\]
In fact $|Y_\theta|$ is bounded above by the greater of the two endpoint
speeds, at most $2\sqrt3<4$, by linear interpolation and the derivative
formulas for the endpoint arclengths. These polynomial bounds turn
inverse-chart evaluation into a bounded grid search. To find parameters
to prescribed tolerance choose an image tolerance small enough by
\eqref{eq:invmod}, enumerate a parameter grid fine enough by the forward
Lipschitz bounds, and use rational approximations to choose a grid image
within that tolerance. The first inequality applies to points named in
spatial coordinates and the second to planar names, including the
planar chords used in lifting paths.
